%%%%%%%%%%%%%%%%%%%%%%%%%%%%%%%%%%%%%%%%%
% baposter Landscape Poster
% LaTeX Template
% Version 1.0 (11/06/13)
%
% baposter Class Created by:
% Brian Amberg (baposter@brian-amberg.de)
%
% This template has been downloaded from:
% http://www.LaTeXTemplates.com
%
% License:
% CC BY-NC-SA 3.0 (http://creativecommons.org/licenses/by-nc-sa/3.0/)
%
%%%%%%%%%%%%%%%%%%%%%%%%%%%%%%%%%%%%%%%%%

%----------------------------------------------------------------------------------------
%	PACKAGES AND OTHER DOCUMENT CONFIGURATIONS
%----------------------------------------------------------------------------------------

\documentclass[landscape,a3paper,fontscale=0.78]{baposter} % Adjust the font scale/size here

%\usepackage[czech]{babel}	
%\usepackage{lmodern}
%\usepackage[T1]{fontenc}
\usepackage[utf8]{inputenc}

\usepackage{graphicx} % Required for including images
\graphicspath{{figures/}} % Directory in which figures are stored

\usepackage{amsmath} % For typesetting math
\usepackage{amssymb} % Adds new symbols to be used in math mode

\DeclareMathOperator{\sgn}{sgn}
\DeclareMathOperator*{\argmax}{argmax}
\DeclareMathOperator*{\softmax}{softmax}
\DeclareMathOperator*{\argsoftmax}{argsoftmax}
\DeclareMathOperator*{\argmin}{argmin}
\DeclareMathOperator{\atan2}{atan2}
\newcommand{\norm}[1]{\lVert #1 \rVert}


\usepackage{booktabs} % Top and bottom rules for tables
\usepackage{enumitem} % Used to reduce itemize/enumerate spacing
\usepackage{palatino} % Use the Palatino font
\usepackage[font=small,labelfont=bf]{caption} % Required for specifying captions to tables and figures



\usepackage{multicol} % Required for multiple columns
\setlength{\columnsep}{1.5em} % Slightly increase the space between columns
\setlength{\columnseprule}{0mm} % No horizontal rule between columns

\usepackage{tikz} % Required for flow chart
\usetikzlibrary{shapes,arrows} % Tikz libraries required for the flow chart in the template

\newcommand{\compresslist}{ % Define a command to reduce spacing within itemize/enumerate environments, this is used right after \begin{itemize} or \begin{enumerate}
\setlength{\itemsep}{1pt}
\setlength{\parskip}{0pt}
\setlength{\parsep}{0pt}
}

\definecolor{lightblue}{rgb}{0.145,0.6666,1} % Defines the color used for content box headers

\begin{document}

\begin{poster}
{
headerborder=closed, % Adds a border around the header of content boxes
colspacing=1em, % Column spacing
bgColorOne=white, % Background color for the gradient on the left side of the poster
bgColorTwo=white, % Background color for the gradient on the right side of the poster
borderColor=lightblue, % Border color
headerColorOne=black, % Background color for the header in the content boxes (left side)
headerColorTwo=lightblue, % Background color for the header in the content boxes (right side)
headerFontColor=white, % Text color for the header text in the content boxes
boxColorOne=white, % Background color of the content boxes
textborder=roundedleft, % Format of the border around content boxes, can be: none, bars, coils, triangles, rectangle, rounded, roundedsmall, roundedright or faded
eyecatcher=true, % Set to false for ignoring the left logo in the title and move the title left
headerheight=0.1\textheight, % Height of the header
headershape=roundedright, % Specify the rounded corner in the content box headers, can be: rectangle, small-rounded, roundedright, roundedleft or rounded
headerfont=\Large\bf\textsc, % Large, bold and sans serif font in the headers of content boxes
%textfont={\setlength{\parindent}{1.5em}}, % Uncomment for paragraph indentation
linewidth=2pt % Width of the border lines around content boxes
}
%----------------------------------------------------------------------------------------
%	TITLE SECTION 
%----------------------------------------------------------------------------------------
%
{%\includegraphics[height=4em]{logo.png}
} % First university/lab logo on the left
{\bf\textsc{Deep Closest Point}\vspace{0.5em}} % Poster title
{\textsc{\{ Nikolay Tsoy, Vojtěch Nydrle \}}} % Author names and institution
{%\includegraphics[height=4em]{logo.png}
} % Second university/lab logo on the right

%----------------------------------------------------------------------------------------
%	ZADANI
%----------------------------------------------------------------------------------------

\headerbox{Zadání}{name=zadani,column=0,row=0}{

\textbf{Tema OJ1:}

Vyzkoušet DCP na real-world datech jako náhrada standardního SLAM algoritmu.

DCP je DL varianta standardního ICP algoritmu. (https://arxiv.org/abs/1905.03304).
%\vspace{0.3em} % When there are two boxes, some whitespace may need to be added if the one on the right has more content
}

%----------------------------------------------------------------------------------------
%	DCP
%----------------------------------------------------------------------------------------

\headerbox{DCP - Deep Closest Point}{name=dcp,column=1,row=0,span=2}{
\begin{multicols}{2}
Na rozdíl od ICP nejprve stanoví přiřezaní odpovídajících si bodů $y=m(\vec x), \vec x \in Y, \vec x \in Y$. Poté z odpovídajících si dvojic $\vec x_i$ a $\vec y_i = m(\vec x_i)$ vypočte rotaci $\textbf{R}_{xy}$ a translaci $\vec{t}_{xy}$ mezi množinami $X$ a $Y$.
Těžiště množin je:
$$ \vec x_c = \frac{1}{N}\sum_{i=1}^N\vec x_i,~\vec y_c = \frac{1}{N}\sum_{i=1}^N\vec y_i$$
a jejich cross-covariantní matice:
   $$ \textbf{H} = \sum_{i=1}^N (\vec x_i-\vec x_c)(\vec y_{m(x_i)}-\vec y_c)^T = \textbf{USV}^T$$
ze které se nechá vypočíst posun a rotace:
   $$\displaystyle \textbf{R}_{xy}=\textbf{VU}^T, ~ \vec{t}_{xy}=\vec y_c - \textbf{R}_{xy}\vec x_c.$$

%\vspace{1.5em}
Pro stanovení odpovídajících si bodů používá PointNet nebo DGCNN skácený o poslední vrstvu, tím ke každému bodu $\vec x_i$ dostane mnohorozměrný vektor $\vec x_i^L$, popisující význam tohoto bodu. Tedy z množin $X$ a $Y$ vypočítá množiny $\textit{F}_X=\lbrace \vec x_1^L, \ldots, \vec x_N^L  \rbrace \subset \mathbb{R}^P $ a $\textit{F}_Y=\lbrace \vec y_1^L, \ldots, \vec y_N^L  \rbrace \subset \mathbb{R}^P$, použitím naučených parametrů $\theta$ použité neuronové sítě. 
  
Pak vypočte přiřazení odpovídajících si bodů: 
  $$m(x_i)=\argsoftmax_{y \in Y}(\Phi_Y(y)^T \Phi_X(x_i)),$$
kde $\Phi$ funkce reagující na společné znaky množin.
$$\Phi_Y(y)=y^L+\phi(y_i^L,\textit{F}_X)$$
$$\Phi_X(x_i)=x_i^L+\phi(x_i^L,\textit{F}_Y)$$
$\phi$ je asymetrická funkce zvaná "Transformer", používaná při zarovnání sekvencí, jejíž parametry se také učí  při učení celého DCP modelu.

Při učení se používá ztrátová funkce
 $$ L=\norm{\textbf{R}_{xy}^T\textbf{R}_{xy}^g-I}^2 + \norm{\vec{t}_{xy}-\vec{t}_{xy}^g}^2+\lambda \norm{\theta},$$
  kde $\textbf{R}_{xy}^g$ a $\vec{t}_{xy}^g$ popisují skutečnou transformaci.
\end{multicols}
}

%----------------------------------------------------------------------------------------
%	RESULTS 1
%----------------------------------------------------------------------------------------

\headerbox{Výsledky autorů}{name=results,column=3,row=0,bottomaligned=dcp}{
Autoři původního článku prováděli experimenty na datasetu ModelNet40, který obsahuje 12 311 CAD modelů rozdělených do 40 kategorii. Použili 9 843 modelů na trénování a 2 468 na testování, modely velikostně přizpůsobily aby byly vepsány jednotkové kouli, náhodně otáčely v rozsahu $45^\circ$ a posouvaly až o $0.5$.

Porovnání DCP a ostatních algoritmu je v tabulce.
\begin{center}
\begin{tabular}{l l l l l l l}
\toprule
Model 				& RMSE(R)		& RMSE(t)	\\
\midrule
ICP					& 29.9148		& 0.2909		\\
Go-ICP 			& 11.8523		& 0.0257		\\
FGR 					& 9.3628			& 0.0139		\\
PointNetLK 	& 15.0954		& 0.0221		\\
\midrule
DCP					& 1.143385		& 0.001786	\\
\bottomrule
\end{tabular}
\captionof{table}{Porovnání s jinými metodami}
\end{center}
Dále ukázaly že s DGCNN dosahuje DCP lepších výsledku než s PointNet, a je možné zpřesnit výsledky následným použitím IPC.
}

%----------------------------------------------------------------------------------------
%	Naše výsledky
%----------------------------------------------------------------------------------------

\headerbox{Dosažené výsledky}{name=vysledky,column=2,span=2,row=0,below=dcp,above=bottom}{

\begin{multicols}{2}

Lorem ipsum im es ....

\end{multicols}
}

%----------------------------------------------------------------------------------------
%	ICP
%----------------------------------------------------------------------------------------

\headerbox{Iterative Closest Point}{name=icp,column=0,below=zadani}{ % This block's bottom aligns with the bottom of the conclusion block
ICP je iterační algoritmus, který určuje rotaci $\textbf{R}_{xy}$ a translaci $\vec{t}_{xy}$ mezi množinami bodů $X$ a $Y$ splňujícími:
\begin{itemize}\compresslist
  	\item $X=\lbrace \vec x_1, \ldots, \vec x_N \rbrace \subset \mathbb{R}^3$
  	\item $Y=\lbrace \vec y_1, \ldots, \vec y_N \rbrace \subset \mathbb{R}^3$
  	\item $Y$ vznikne otočením a posunutím $X$
\end{itemize}
tak že minimalizuje chybu $E(\textbf{R}_{xy},\vec{t}_{xy})$.
$$ E(\textbf{R}_{xy},\vec{t}_{xy})=\frac{1}{N}\sum_{i=1}^N \norm{\textbf{R}_{xy}\vec x_i+\vec{t}_{xy} -\vec y_{m(x_i)}}^2 $$
$$\displaystyle m(x_i)=\argmin_j \norm{\textbf{R}_{xy}\vec x_i+\vec{t}_{xy} -\vec y_j}^2 $$
 \textbf{Problémy IPC}
  \begin{enumerate}\compresslist
  \item Nelze optimalizovat  $\textbf{R}_{xy}$, $\vec{t}_{xy}$ i $m$ najednou proto v jednom kroku optimalizuje $m$ a v dalším  $\textbf{R}_{xy}$ a $\vec{t}_{xy}$ 
  \item velmi často najde jen lokální optimum
  \item neuvažuje zajímavost některých bodů 
  \item neporadí si se šumem a řídkostí měření 
  \end{enumerate}  
}


%----------------------------------------------------------------------------------------
%	DALŠÍ VÝZKUM	
%----------------------------------------------------------------------------------------

\headerbox{Další možné zlepšení}{name=next,column=1,below=dcp,bottomaligned=vysledky}{
Dalšího zlepšení, které jsme již nestihly vyzkoušet, by šlo dosáhnout následujícími způsoby:
\begin{itemize}\compresslist
\item Natrénovat neuronovou sít na všech možných kombinacích záběrů, stejné scény, ne jen na dvou po sobě jdoucích.
\item Trénovat neuronovou sít na více jízdách než jen na jedné.
\item Při výpočtu aktuální pozice nevycházet jen z předešlého záběru scény ale všech v dosahu lidaru.
\item Po vypočtení pozice pomocí DCP použít na tento výsledek ještě ICP pro zpřesnění výsledku.
\end{itemize}

}
%----------------------------------------------------------------------------------------
%	DATASET
%----------------------------------------------------------------------------------------

\headerbox{Dataset}{name=datasetcolumn=0, below=icp, above=bottom}{ % This block's bottom aligns with the bottom of the conclusion block

Jako real-world data byl vybrán "KITTI Vision" dataset, obsahující 22 nahrávek jízd autem s lidarem GPS a spoustou dalších senzorů. K prvním 11 jízdám jsou poskytnuta jak lidarová data, ta skutečná pozice a natočení auta.

DCP bylo trénováno na jedné jízdě a testováno na jiné.


}

%----------------------------------------------------------------------------------------

\end{poster}

\end{document}